% RESEARCH_DRAFT_SHA256: d16c0de080f27f7826739583df21ef921f7c3fd7b647f96dd2eb2502f4290cae
% Layout and source-reference wrappers expanded for the web renderer.
\documentclass[11pt]{article}
\usepackage[margin=25mm]{geometry}
\usepackage{amsthm}
\usepackage{fontspec,unicode-math}
\setmainfont{Latin Modern Roman}
\setmathfont{Latin Modern Math}
\usepackage{amsmath,amssymb,mathtools}
\usepackage{xurl,hyperref}
\hypersetup{colorlinks=true,urlcolor=blue,linkcolor=blue}
\setcounter{secnumdepth}{0}
\setlength{\parindent}{0pt}
\setlength{\parskip}{5pt}
\setlength{\emergencystretch}{3em}
\newcommand{\sref}[2]{\hyperlink{source:#1:#2}{[S#2]}}
\newcommand{\cataloguescope}[1]{\paragraph{Source question.}#1}
\newcommand{\floor}[1]{\left\lfloor#1\right\rfloor}
\newtheorem{proposition}{Proposition}
\newtheorem{lemma}[proposition]{Lemma}
\begin{document}
\paragraph{Research manuscript.} OpenAI \texttt{gpt-6-astra}, reasoning \texttt{ultra}; 11 October 2026.
\paragraph{Preserved source body.} The complete catalogue mathematical text and sources follow with layout and source links expanded, including source qualifications and its historical preparation notes. New research begins at the section headed Research attempt.

% UnsolvedMath ID 30002627; source catalogue code OWR-13105-006
\section{Witness Pseudorandom Functions from Witness Encryption}\label{30002627}
{\small UnsolvedMath ID 30002627 · Computer Science}
\subsection{Definitions and mathematical statement}

\paragraph{Shared notation.}Standard number systems are denoted by $\mathbb N,\mathbb Z,\mathbb Q,\mathbb R,\mathbb C$; index ranges, computational models and specialized notation follow the cited problem source and the definitions below.


Witness encryption encrypts to an instance of an NP relation and permits decryption using a valid witness. A witness pseudorandom function associates pseudorandom outputs with a relation and gives an appropriate witness-holder evaluation capability.

A cryptographic security requirement compares an efficient adversary\textquotesingle s advantage with a negligible function of the security parameter. A negligible function is smaller than every inverse polynomial eventually; computational and information-theoretic security are different requirements.
\paragraph{Statement recorded in the supplied source.}
Can witness pseudorandom functions be constructed generically from any witness encryption scheme? \href{https://doi.org/10.4171/owr/2014/35}{[S2]}
\paragraph{Scope qualification.}
The security notions and meaning of generic construction are report-specific.
\paragraph{Source status.}
The supplied catalogue labels this record open. Its source-status metadata was checked against the supplied snapshot on 2026-10-11. This import does not independently establish current open status or mathematical correctness. The catalogue literature assessment is dated 2026-08-21.
\subsection{Short English statement}
Construct witness pseudorandom functions generically from arbitrary witness-encryption schemes.
\subsection{Sources}
\begin{itemize}
\item[S1] Ulam AI and the UnsolvedMath Contributors, UnsolvedMath: A Curated Collection of Open Mathematics Problems, record 30002627 (OWR-13105-006). Catalogue curation and metadata: CC BY 4.0. The supplied source question is retained; definitions and the English summary are editorial additions.
\url{https://huggingface.co/datasets/ulamai/UnsolvedMath}.
\item[S2] Cryptography (2015), Oberwolfach Reports, DOI 10.4171/owr/2014/35, pp. 1933--1990
\url{https://doi.org/10.4171/owr/2014/35}.
\item[S3] Zhandry witness PRF paper
\url{https://eprint.iacr.org/2014/301}.
\item[S4] TCC 2016 witness PRFs
\url{https://doi.org/10.1007/978-3-662-49099-0_16}.
\textit{Underlying Oberwolfach Reports material retains its source-specific licence, including CC BY-SA 4.0 where applicable; see THIRD\_PARTY\_NOTICES.md.}
\end{itemize}
\paragraph{Preparation and review.}
Original question attributed to the cited source and UnsolvedMath contributors. Codex (GPT-6; exact serving revision not exposed) assisted with definitions, English summaries, source inspection and formatting. No new human mathematical authorship or independent proof verification is asserted. Canonical ID 30002627 and source code OWR-13105-006 checked against the complete supplied registry and existing repository statements; known flagged repetitions and obvious identical conjectures excluded. Full mathematical-equivalence/literature review has not been performed. Statement preparation and final review followed contribution prompts C and D; automated checks verify structure and identity, not open status.
\label{end:30002627}

\section{Research attempt}
\subsection{Literature, source context and scope}
The full original catalogue body and source qualifications are preserved above. The relevant primary workshop passages were read in the cached publisher report. Historical claims and subsequent attributed results are distinguished from the mathematical work below. A restricted example, standard lemma or known theorem is not counted as new progress on the entire original target. Missing assumptions and repeated or rhetorical extracts are not silently turned into different solved problems.

\subsection{A finite-domain witness-encryption table, and its size obstruction}
The report's witness-PRF must evaluate publicly with a valid witness while keeping its value pseudorandom at false instances even given its public evaluation key and permitted function queries. Its known construction uses multilinear maps; the requested generic witness-encryption implication is broader. We test an exhaustive-table construction under ordinary witness-encryption correctness and false-instance message indistinguishability.
Let $X$ be an explicitly listed domain of size $D$ polynomial in the security parameter, let $R(x,w)$ be the witness relation, and choose independent uniform $\ell$-bit strings $r_x$. Store their table in a secret key. Define $F(x)=r_x$ and publish $c_x=\operatorname{WE.Enc}(x,r_x)$ for every $x\in X$. On input $(x,w)$, first verify $R(x,w)$, returning $\bot$ if invalid; otherwise decrypt $c_x$ with $w$.
\begin{proposition}
This construction has witness-evaluation correctness. For a fixed false challenge instance $x_*\in X$, its value is computationally indistinguishable from an independent uniform string, even with the public table and adaptive value queries at points other than $x_*$, assuming false-instance security of witness encryption. This is a selective, single-challenge assertion.
\end{proposition}
\begin{proof}
On a valid witness, witness-encryption correctness returns $r_x=F(x)$; invalid witnesses are rejected explicitly. For security, choose independent uniform strings $u,v$ and submit them as the two messages in a witness-encryption challenge for the fixed false instance $x_*$. Insert the challenge ciphertext as $c_{x_*}$, generate every other table entry honestly, and give $u$ as the function-value challenge. Queries at other points can be answered from their stored strings. If the ciphertext encrypts $u$, the joint distribution is exactly that of the true function value at $x_*$. If it encrypts $v$, the public table encrypts an independent uniform true value while the announced challenge $u$ is independent uniform. All other table entries and query answers agree in distribution. Thus any distinguisher for these two experiments gives one with the same advantage against witness encryption. The simulation is polynomial because $D$ is polynomial and all other operations are efficient by the primitive's hypotheses.
\end{proof}
\subsection{Verification and remaining gaps}
The independent-message hybrid was checked by enumerating all 16 ordered pairs of two-bit strings, including collisions. Table storage for the full $n$-bit domain has at least $2^n\ell$ secret bits, and that growth was evaluated for several $n$. The failed step in extending this construction is not correctness but efficiency: enumerating the entire ordinary domain $\{0,1\}^n$ is exponential in $n$. The proof also supplies only the stated selective single-point security game, not a proof for an unrestricted adaptively chosen challenge. No compression of the table, generic full-domain implication, or weaker-assumption witness-PRF construction is established. The 2014 primary construction is attributed, not independently reproduced.

\section{References and source attribution}
Primary workshop report: \url{https://doi.org/10.4171/owr/2014/35}. Relevant passages read: 2014/35 pp. 1974--1975, witness-PRF syntax, false-instance pseudorandomness, oracle queries and the witness-encryption question.

\url{https://eprint.iacr.org/2014/301}. Read primary Zhandry abstract defining witness PRFs and presenting a multilinear-map construction; historical security descriptions are not asserted as current guarantees.

All retrievals were checked on 11 October 2026. Original question attribution remains in the preserved Sources section. External theorems are attributed and not claimed as new discoveries or fully verified merely by citation.


\paragraph{Integration review.} Codex (GPT-6; exact serving revision not exposed) reviewed the written arguments and expanded layout and source-reference wrappers. Repository metadata and the production web renderer were checked using Python 3.13 and JavaScript with MathJax 3.2.2. This model-assisted review is not independent human proof verification; no PDF compilation or formal proof checking is asserted.

\section{Completion Estimate}
\noindent\textbf{COMPLETION: 0\%}
No substantiated new progress toward the entire original target is claimed. The worked results and failed route are delimited above; a known result, illustrative model, or auxiliary obstruction is not counted as original-target completion.

\section{Final status and attribution}
\textbf{Attempt status: unresolved.} This describes this research attempt, not a new ruling on the source catalogue's status or on external literature claims.
The mathematics and this self-audit were generated by OpenAI \texttt{gpt-6-astra}, reasoning \texttt{ultra}, on 11 October 2026. Codex (GPT-6; exact serving revision not exposed) orchestrated the batch, reviewed the written mathematical arguments, and contributed formatting/repository review. No human mathematical authorship, independent human proof checking, or formal verification is claimed. The historical preparation paragraphs in the reproduced catalogue body concern that body only.

Contribution prompt A was applied to the complete definition and supplied source record, followed by the adversarial checks of prompt B. Local reproduction files are \texttt{batch019/checks.py}, \texttt{computations.json}, and the per-ID audit JSON. They were not public links at preparation time. Numerical checks are finite and do not prove unrestricted results. No LaTeX engine was found locally; compilation was not performed by this research agent.

\end{document}
